\section{The Schmidt number and the full reduction criterion}

The Schmidt number of a bipartite state is the smallest $n$ for which the state has
a convex decomposition into pure states each of Schmidt rank at most $n$. Separable
states are precisely those of Schmidt number one, and a state of Schmidt number at
most $n$ satisfies the reduction criterion
\begin{align}
    n\,\rho_1\otimes\Id & \ge\rho,
    \notag                         \\
    n\,\Id\otimes\rho_2 & \ge\rho,
    \notag
\end{align}
which is \cite[Equation~(3.18)]{Wolf2012Quantum}. The bound on the Schmidt number is
recorded as a finite sum of pure-state projectors of Schmidt rank at most $n$, the
unnormalized form of the convex decomposition; rescaling a vector by the square root
of its weight scales its coefficient matrix and so leaves the Schmidt rank unchanged.

\begin{definition}[Bounded Schmidt number]
    \label{def:has_schmidt_number_le}
    \lean{Matrix.HasSchmidtNumberLE}
    \leanok
    \uses{def:schmidt_rank}
    A bipartite matrix $\rho$ on $\MN{d}\otimes\MN{d'}$ has Schmidt number at most $n$
    when it is a finite sum of pure-state projectors of Schmidt rank at most $n$,
    \begin{align}
        \rho        & =\sum_i\ket{\psi_i}\bra{\psi_i},
        \notag                                         \\
        \SR(\psi_i) & \le n.
        \notag
    \end{align}
\end{definition}

\begin{theorem}[Basic properties of bounded Schmidt number]
    \label{thm:has_schmidt_number_le_basic}
    \lean{Matrix.HasSchmidtNumberLE.posSemidef, Matrix.HasSchmidtNumberLE.add,
        Matrix.HasSchmidtNumberLE.mono, Matrix.HasSchmidtNumberLE.smul}
    \leanok
    \uses{def:has_schmidt_number_le}
    A matrix of Schmidt number at most $n$ is positive semidefinite. States of
    Schmidt number at most $n$ are closed under addition and under multiplication by a
    non-negative real scalar, and a bound on the Schmidt number relaxes to any larger
    bound.
\end{theorem}

\begin{proof}
    \leanok
    Each summand $\ket{\psi_i}\bra{\psi_i}$ is a rank-one positive semidefinite
    matrix, and a finite sum of positive semidefinite matrices is positive
    semidefinite. Concatenating two pure-state families realizes a sum as a single
    finite sum of pure-state projectors of Schmidt rank at most $n$, and a pure
    summand of Schmidt rank at most $n$ also has Schmidt rank at most any larger
    bound. Multiplication by a non-negative real $a$ rescales each pure summand
    $\psi_i$ to $\sqrt a\,\psi_i$, which leaves its Schmidt rank unchanged and
    reproduces $a$ times the original state.
\end{proof}

\begin{theorem}[Schmidt number one is separability]
    \label{thm:has_schmidt_number_le_one_iff_separable}
    \lean{Matrix.hasSchmidtNumberLE_one_iff_isSeparable}
    \leanok
    \uses{def:has_schmidt_number_le, def:is_separable, def:schmidt_rank,
        lem:matrix_rank_factorization_coordinate, thm:schmidt_rank_basic_bounds,
        thm:is_separable_kronecker, thm:is_separable_add}
    A bipartite matrix has Schmidt number at most one if and only if it is separable.
\end{theorem}

\begin{proof}
    \leanok
    A pure state of Schmidt rank at most one has a coefficient matrix of rank at most
    one, which factors as $\psi(i,j)=u_i v_j$; its projector is the Kronecker product
    $\ket{u}\bra{u}\otimes\ket{v}\bra{v}$ of two rank-one positive semidefinite
    matrices, hence separable, and separable matrices are closed under addition.
    Conversely a Kronecker product $A\otimes B$ of positive semidefinite matrices,
    after expanding each factor into rank-one projectors, is a double sum of
    product-vector projectors, each of Schmidt rank at most one, so it has Schmidt
    number at most one; a separable matrix is a finite sum of such products.
\end{proof}

\begin{theorem}[Convexity of the Schmidt-number set]
    \label{thm:convex_setOf_has_schmidt_number_le}
    \lean{Matrix.convex_setOf_hasSchmidtNumberLE}
    \leanok
    \uses{def:has_schmidt_number_le}
    For each $n$ the set $S_n$ of bipartite states of Schmidt number at most $n$ is
    convex.
\end{theorem}

\begin{proof}
    \leanok
    A convex combination $a x + b y$ with $a,b\ge 0$ is a sum of two non-negatively
    scaled states of Schmidt number at most $n$, and that bound is preserved under
    non-negative scaling and under addition; rescaling a vector by a non-negative
    scalar scales its coefficient matrix and so does not increase the Schmidt rank.
\end{proof}

\begin{theorem}[Compactness of the Schmidt-number state set]
    \label{thm:compact_setOf_has_schmidt_number_le_trace_one}
    \lean{Matrix.isCompact_setOf_hasSchmidtNumberLE_trace_one}
    \leanok
    \uses{def:has_schmidt_number_le, def:schmidt_rank,
        thm:convex_setOf_has_schmidt_number_le}
    For each $n$ the set $S_n$ of trace-one bipartite states of Schmidt number at most
    $n$ is compact.
\end{theorem}

\begin{proof}
    \leanok
    A trace-one state of Schmidt number at most $n$ is a convex combination of pure
    states $\ket{\psi}\bra{\psi}$ whose vector $\psi$ has unit norm and Schmidt rank
    at most $n$, so $S_n$ is the convex hull of the set $P_n$ of those pure states.
    The set $P_n$ is the image, under the continuous projector map
    $\psi\mapsto\ket{\psi}\bra{\psi}$, of the unit vectors of Schmidt rank at most
    $n$. In the Euclidean norm those unit vectors form the unit sphere, which is
    compact in finite dimension, intersected with the closed set on which the Schmidt
    rank is at most $n$ (the coefficient matrix depends continuously on the vector and
    the matrices of rank at most $n$ form a closed set); hence $P_n$ is compact. The
    trace of $\ket{\psi}\bra{\psi}$ is the squared Euclidean norm of $\psi$, so the
    trace-one constraint is exactly the unit-norm constraint and the weights of the
    convex decomposition are the squared norms, summing to the trace. The convex hull
    of a compact set is compact in a finite-dimensional real normed space.
\end{proof}

\begin{theorem}[Pure-state step of the reduction criterion]
    \label{thm:tensor_map_id_t_eta_posSemidef_pure}
    \lean{Matrix.tensorMapId_tEta_posSemidef_of_hasSchmidtRankLE}
    \leanok
    \uses{def:schmidt_rank, def:t_eta, def:tensor_map_id, thm:t_eta_threshold,
        thm:square_schmidt_rank_exact_parametrization,
        lem:rank_one_ampliation_choi_compression,
        thm:choi_right_compression_quadratic_form, thm:schmidt_rank_choi_test_iff}
    For a pure state $\ket{\psi}\bra{\psi}$ on the square bipartite system
    $\MN{D}\otimes\MN{D}$ with $\SR(\psi)\le n$ and
    $1\le n<D$, applying $T_n$ to the first factor keeps the result positive
    semidefinite: $(T_n\otimes\id)(\ket{\psi}\bra{\psi})\ge 0$.
\end{theorem}

\begin{proof}
    \leanok
    The map $T_n$ is $n$-positive for $1\le n<D$. Write $\psi$ through the maximally
    entangled vector as $\psi=\psi_X$ for a square matrix $X$ on the right factor with
    $\rank(X)=\SR(\psi)\le n$. Then
    $(T_n\otimes\id)(\ket{\psi}\bra{\psi})$ is the right-factor Choi compression of
    $T_n$ by $X$. Its quadratic form on any vector $\eta$ is the Choi quadratic form
    of $T_n$ on $R_X^\dagger\eta$, whose Schmidt rank is at most
    $\rank(X)\le n$. The Schmidt-rank Choi criterion makes that form
    non-negative, so the compression is positive semidefinite.
\end{proof}

\begin{theorem}[Pure-state step for a general positive map]
    \label{thm:tensor_map_id_posSemidef_pure}
    \lean{Matrix.tensorMapId_posSemidef_of_hasSchmidtRankLE}
    \leanok
    \uses{def:schmidt_rank, def:tensor_map_id, def:k_positive_map,
        def:choi_right_compression, thm:square_schmidt_rank_exact_parametrization,
        lem:rank_one_ampliation_choi_compression,
        thm:choi_right_compression_quadratic_form, thm:schmidt_rank_choi_test_iff}
    For a pure state $\ket{\psi}\bra{\psi}$ on the square bipartite system
    $\MN{D}\otimes\MN{D}$ with $\SR(\psi)\le n$ and any $n$-positive map
    $T$, applying $T$ to the first factor keeps the result positive semidefinite:
    $(T\otimes\id)(\ket{\psi}\bra{\psi})\ge 0$.
\end{theorem}

\begin{proof}
    \leanok
    Write $\psi$ through the maximally entangled vector as $\psi=\psi_X$ for a square
    matrix $X$ on the right factor with
    $\rank(X)=\SR(\psi)\le n$. Then
    $(T\otimes\id)(\ket{\psi}\bra{\psi})$ is the right-factor Choi compression of $T$
    by $X$, and its quadratic form on a vector $\eta$ equals the Choi quadratic form
    $C_T$ of $T$ evaluated on $R_X^\dagger\eta$, a vector of Schmidt rank at most
    $\rank(X)\le n$. Because $T$ is $n$-positive,
    $\langle R_X^\dagger\eta,\, C_T\, R_X^\dagger\eta\rangle\ge 0$, so the
    compression is positive semidefinite.
\end{proof}

\begin{theorem}[Positive maps and entanglement, only-if direction]
    \label{thm:has_schmidt_number_le_tensor_map_id}
    \lean{Matrix.HasSchmidtNumberLE.tensorMapId_posSemidef}
    \leanok
    \uses{def:has_schmidt_number_le, def:tensor_map_id, def:k_positive_map,
        thm:tensor_map_id_sum, thm:tensor_map_id_posSemidef_pure}
    A bipartite state on the square system $\MN{D}\otimes\MN{D}$ of Schmidt number at
    most $n$ satisfies $(T\otimes\id)(\rho)\ge 0$ for every $n$-positive map $T$.
    This is the square case $d=d'=D$ of the only-if direction of Wolf's
    Proposition~3.4 for a general bipartite system
    $\MN{d}\otimes\MN{d'}$ \cite[Chapter~3, Proposition~3.4]{Wolf2012Quantum}.
\end{theorem}

\begin{proof}
    \leanok
    The state is a finite sum of pure-state projectors of Schmidt rank at most $n$.
    Applying $T$ to the first factor is linear, so it distributes over the sum; the
    pure-state step makes each summand positive semidefinite, and a finite sum of
    positive semidefinite matrices is positive semidefinite.
\end{proof}

\begin{theorem}[Only-if direction for a general bipartite system]
    \label{thm:has_schmidt_number_le_tensor_map_id_general}
    \lean{Matrix.HasSchmidtNumberLE.tensorMapId_posSemidef_general}
    \leanok
    \uses{def:has_schmidt_number_le, def:tensor_map_id, def:k_positive_map,
        thm:has_schmidt_number_le_tensor_map_id}
    On a general bipartite system $\MN{d}\otimes\MN{d'}$, with no relation imposed
    between the two factors, a state of Schmidt number at most $n$ satisfies
    $(T\otimes\id)(\rho)\ge 0$ for every $n$-positive map $T$ on $\MN{d}$.
    This is the only-if direction of Wolf's Proposition~3.4, the first step of
    eq.~(3.18) \cite[Chapter~3, Proposition~3.4]{Wolf2012Quantum}, with no restriction
    on the bipartite dimensions.
\end{theorem}

\begin{proof}
    \leanok
    Pad the bipartite state to the square system $\MN{D}\otimes\MN{D}$ with
    $D=\max(d,d')$ by zero rows and columns: when $d'\le d$ the second factor is padded
    up to $d$; when $d\le d'$ the first factor is padded up to $d'$ and $T$ is extended
    to a map on $\MN{d'}$ that preserves $n$-positivity. Padding preserves the
    Schmidt-number bound, the square result makes the padded ampliation positive
    semidefinite, and the original ampliation is its principal submatrix, hence positive
    semidefinite.
\end{proof}

\begin{theorem}[Reduction step from the Schmidt-number premise]
    \label{thm:has_schmidt_number_le_tensor_map_id_t_eta}
    \lean{Matrix.HasSchmidtNumberLE.tensorMapId_tEta_posSemidef}
    \leanok
    \uses{def:has_schmidt_number_le, def:t_eta, def:tensor_map_id,
        thm:tensor_map_id_sum, thm:tensor_map_id_t_eta_posSemidef_pure}
    For $1\le n<D$, a bipartite state on the square system
    $\MN{D}\otimes\MN{D}$ of Schmidt number at most $n$ satisfies
    $(T_n\otimes\id)(\rho)\ge 0$.
\end{theorem}

\begin{proof}
    \leanok
    The state is a finite sum of pure-state projectors of Schmidt rank at most $n$.
    Applying $T_n$ to the first factor is linear, so it distributes over the sum; the
    pure-state step makes each summand positive semidefinite, and a finite sum of
    positive semidefinite matrices is positive semidefinite.
\end{proof}

\begin{theorem}[Reduction criterion with the Schmidt-number premise]
    \label{thm:reduction_criterion_of_schmidt_number}
    \lean{Matrix.reductionCriterion_left_of_hasSchmidtNumberLE_general,
        Matrix.reductionCriterion_right_of_hasSchmidtNumberLE_general}
    \leanok
    \uses{def:has_schmidt_number_le, def:t_eta,
        thm:has_schmidt_number_le_tensor_map_id_general, thm:reduction_criterion}
    On a general bipartite system $\MN{d}\otimes\MN{d'}$, a state of Schmidt number at
    most $n$ satisfies the reduction criterion
    \begin{align}
        n\,\Id\otimes\rho_2 & \ge\rho
                            &         & (1\le n<d),
        \notag                                       \\
        n\,\rho_1\otimes\Id & \ge\rho
                            &         & (1\le n<d'),
        \notag
    \end{align}
    with $\rho_1$ and $\rho_2$ the reduced densities on the first and second factors.
    No relation is imposed between $d$ and $d'$; each inequality carries only the
    $n$-positivity threshold of $T_n$ on its acting factor.
\end{theorem}

\begin{proof}
    \leanok
    The Schmidt-number premise gives $(T_n\otimes\id)(\rho)\ge 0$ on the general system
    (first step), and the dimension-general operator implication then yields
    $n\,\Id\otimes\rho_2\ge\rho$. The factor swap of a state of Schmidt number at most
    $n$ again has Schmidt number at most $n$, since it sends each pure summand $\psi$ to
    $\psi\circ\mathrm{swap}$, whose Schmidt coefficient matrix is the transpose
    $M_{\psi\circ\mathrm{swap}}=M_\psi^{\mathsf T}$ of that of $\psi$ and so has the same
    rank; applying the first step to the swapped state (whose first factor is $d'$) and
    the operator implication in its symmetric form gives $n\,\rho_1\otimes\Id\ge\rho$.
\end{proof}


\begin{theorem}[Trace-form representation of a real-linear functional]
    \label{thm:exists_isHermitian_trace_form_re}
    \lean{Matrix.exists_isHermitian_trace_form_re}
    \leanok
    Every real-linear functional $g$ on the space of square complex matrices is the
    trace form of a Hermitian matrix: there is a Hermitian $H_0$ with
    $g(X)=\re\tr(X H_0)$ for every Hermitian $X$.
\end{theorem}

\begin{proof}
    \leanok
    The real bilinear pairing $(X,H)\mapsto\re\tr(X H)$ is nondegenerate,
    because on the diagonal pair $(H,H^{\dagger})$ it equals the squared Frobenius norm
    $\re\tr(H H^{\dagger})=\sum_{i,j}\lvert H_{ij}\rvert^2$, which is positive
    unless $H=0$. Hence the map sending $H$ to the functional
    $X\mapsto\re\tr(X H)$ is an injective real-linear map from the matrix
    space to its real dual. These two spaces have the same finite real dimension, so the
    map is also surjective, and every $g$ is represented by some $H$. Replacing $H$ by
    its Hermitian part $\tfrac12(H+H^{\dagger})$ keeps the value on Hermitian inputs,
    since $\tr(X H^{\dagger})=\overline{\tr(X H)}$ when $X$ is Hermitian, and produces a
    Hermitian representing matrix.
\end{proof}

\begin{theorem}[Entanglement witness from a separating hyperplane]
    \label{thm:exists_entanglement_witness}
    \lean{Matrix.exists_isHermitian_witness}
    \leanok
    \uses{def:has_schmidt_number_le, def:schmidt_rank,
        thm:convex_setOf_has_schmidt_number_le,
        thm:compact_setOf_has_schmidt_number_le_trace_one,
        thm:exists_isHermitian_trace_form_re}
    Let $\rho$ be a trace-one Hermitian bipartite state whose Schmidt number exceeds $n$.
    Then there is a Hermitian operator $W$ such that, for every $\psi$ of Schmidt rank
    at most $n$,
    \begin{align}
        \re\tr(W\rho)            & <0,
        \notag                            \\
        \re\bra{\psi}W\ket{\psi} & \ge 0.
        \notag
    \end{align}
    This is the only-if direction of Wolf's Proposition~3.3
    \cite[Chapter~3, Proposition~3.3]{Wolf2012Quantum}: a state of Schmidt number larger
    than $n$ is detected by an entanglement witness for $S_n$.
\end{theorem}

\begin{proof}
    \leanok
    The set $S_n$ of trace-one states of Schmidt number at most $n$ is convex and
    compact, and $\rho\notin S_n$. The geometric Hahn--Banach theorem separates the
    closed convex set $\{\rho\}$ from the compact convex set $S_n$ by a continuous
    real-linear functional $f$ and a constant $c$ with $f(\rho)<c<f(\sigma)$ for every
    $\sigma\in S_n$. Representing $f$ as the trace form of a Hermitian $H_0$ and setting
    $W=H_0-c\,\Id$ gives $\re\tr(W\rho)=f(\rho)-c<0$. For a vector $\psi$
    of Schmidt rank at most $n$ the projector $\ket{\psi}\bra{\psi}$, after
    normalization by $\lVert\psi\rVert^2$, is a trace-one state of $S_n$, so its value
    under $f$ is at least $c$; multiplying back by the squared norm gives
    $\re\bra{\psi}W\ket{\psi}\ge 0$, with the zero vector giving the value
    $0$ directly.
\end{proof}

\begin{theorem}[Witness criterion for Schmidt number]
    \label{thm:not_has_schmidt_number_le_iff_exists_witness}
    \lean{Matrix.not_hasSchmidtNumberLE_iff_exists_witness}
    \leanok
    \uses{def:has_schmidt_number_le, def:schmidt_rank,
        thm:exists_entanglement_witness}
    Let $\rho$ be a trace-one Hermitian bipartite state. Then the Schmidt number of
    $\rho$ exceeds $n$ if and only if there is a Hermitian operator $W$ such that, for
    every $\psi$ of Schmidt rank at most $n$,
    \begin{align}
        \re\tr(W\rho)            & <0,
        \notag                            \\
        \re\bra{\psi}W\ket{\psi} & \ge 0.
        \notag
    \end{align}
    This is Wolf's Proposition~3.3 \cite[Chapter~3, Proposition~3.3]{Wolf2012Quantum}: a
    state has Schmidt number larger than $n$ exactly when it is detected by an
    entanglement witness for $S_n$.
\end{theorem}

\begin{proof}
    \leanok
    The forward implication is the separating-hyperplane construction
    (\ref{thm:exists_entanglement_witness}). For the converse, suppose $\rho$ had Schmidt
    number at most $n$, so $\rho=\sum_i\ket{\psi_i}\bra{\psi_i}$ with each $\psi_i$ of
    Schmidt rank at most $n$. By linearity of the trace,
    $\re\tr(W\rho)=\sum_i\re\tr\!(W\ket{\psi_i}\bra{\psi_i})$,
    a sum of non-negative terms, contradicting $\re\tr(W\rho)<0$. This
    converse uses no density-matrix hypotheses on $\rho$.
\end{proof}

\begin{theorem}[Surjectivity of the Choi correspondence]
    \label{thm:exists_choiMatrix_eq}
    \lean{ChoiJamiolkowski.exists_choiMatrix_eq}
    \leanok
    \uses{thm:cp_iff_choi_posSemidef}
    Assume $D>0$. Every bipartite matrix $W$ on $\C^D\otimes\C^D$ is the Choi matrix
    $\tau_T=(T\otimes\Id)(\ket{\Omega}\bra{\Omega})$ of some linear map
    $T\colon\MN{D}\to\MN{D}$.
\end{theorem}

\begin{proof}
    \leanok
    The Choi correspondence $T\mapsto\tau_T$ is complex-linear and injective. Its domain,
    the space of linear maps on $\MN{D}$, and its codomain, the bipartite matrix space on
    $\C^D\otimes\C^D$, both have complex dimension $D^4$, so the injective map is also
    surjective. Hence every $W$ is a Choi matrix.
\end{proof}

\begin{theorem}[A witness is the Choi matrix of an $n$-positive map]
    \label{thm:exists_isNPositiveMap_choiMatrix_eq}
    \lean{ChoiJamiolkowski.exists_isNPositiveMap_choiMatrix_eq_of_witness}
    \leanok
    \uses{def:k_positive_map, def:schmidt_rank, thm:exists_choiMatrix_eq,
        thm:schmidt_rank_choi_test_iff, thm:exists_entanglement_witness}
    Assume $D>0$. Let $W$ be a Hermitian operator on $\C^D\otimes\C^D$ with
    $\re\bra{\psi}W\ket{\psi}\ge 0$ for every $\psi$ of Schmidt rank at most $n$.
    Then $W$ is the Choi matrix of an $n$-positive map $T$. This is the
    Choi--Jamio\l kowski translation of the entanglement witness of Wolf's Proposition~3.3
    into the $n$-positive map of Wolf's Proposition~3.4
    \cite[Chapter~3, Proposition~3.4]{Wolf2012Quantum}.
\end{theorem}

\begin{proof}
    \leanok
    By the surjectivity of the Choi correspondence (\ref{thm:exists_choiMatrix_eq}) there
    is a linear map $T$ with $\tau_T=W$. Because $W$ is Hermitian, the Choi quadratic form
    $\bra{\psi}W\ket{\psi}$ is real, and it equals the witness expectation
    $\tr(W\ket{\psi}\bra{\psi})$; the witness condition makes it non-negative on every
    vector of Schmidt rank at most $n$. This is exactly the Schmidt-rank Choi criterion
    (\ref{thm:schmidt_rank_choi_test_iff}) for $n$-positivity of $T$.
\end{proof}

\begin{theorem}[Choi trace pairing through the trace-pairing adjoint]
    \label{thm:trace_choiMatrix_omega_quadratic_adjoint}
    \lean{ChoiJamiolkowski.trace_choiMatrix_mul_eq_omegaVec_quadraticForm_traceAdjointMap}
    \leanok
    \uses{def:trace_pairing_adjoint, thm:trace_pairing_adjoint_ampliation}
    For a linear map $T\colon\MN{D}\to\MN{D}$ with Choi matrix $\tau_T$ and any bipartite
    matrix $\rho$ on $\C^D\otimes\C^D$,
    \begin{align}
        \tr(\tau_T\,\rho)
         & =\bra{\Omega}(T^{*}\otimes\Id)(\rho)\ket{\Omega},
        \notag
    \end{align}
    where $T^{*}$ is the trace-pairing adjoint of $T$.
\end{theorem}

\begin{proof}
    \leanok
    Since $\tau_T=(T\otimes\Id)(\ket{\Omega}\bra{\Omega})$, cyclicity of the trace gives
    $\tr(\tau_T\rho)=\tr(\rho\,(T\otimes\Id)(\ket{\Omega}\bra{\Omega}))$. The
    $\Id$-ampliation of the trace-pairing adjoint is the trace-pairing adjoint of the
    ampliation, so moving $T$ across the trace pairing replaces
    $(T\otimes\Id)$ by $(T^{*}\otimes\Id)$ acting on $\rho$ and turns the pairing against
    $\ket{\Omega}\bra{\Omega}$ into the quadratic form $\bra{\Omega} \cdot\,\ket{\Omega}$.
\end{proof}

\begin{theorem}[Positive maps detect high Schmidt number]
    \label{thm:exists_isNPositiveMap_detects}
    \lean{Matrix.exists_isNPositiveMap_tensorMapId_not_posSemidef}
    \leanok
    \uses{def:k_positive_map, def:has_schmidt_number_le,
        thm:exists_entanglement_witness, thm:exists_isNPositiveMap_choiMatrix_eq,
        thm:k_positive_trace_pairing_adjoint, thm:trace_choiMatrix_omega_quadratic_adjoint}
    Let $\rho$ be a trace-one Hermitian bipartite state on $\C^D\otimes\C^D$ whose Schmidt
    number exceeds $n$. Then there is an $n$-positive map $T\colon\MN{D}\to\MN{D}$ such that
    $(T\otimes\Id)(\rho)$ is not positive semidefinite. This is the if direction of Wolf's
    Proposition~3.4 \cite[Chapter~3, Proposition~3.4]{Wolf2012Quantum}.
\end{theorem}

\begin{proof}
    \leanok
    The entanglement witness $W$ for $\rho$ (\ref{thm:exists_entanglement_witness}) is the
    Choi matrix of an $n$-positive map $P$ (\ref{thm:exists_isNPositiveMap_choiMatrix_eq}).
    Its trace-pairing adjoint $T=P^{*}$ is again $n$-positive
    (\ref{thm:k_positive_trace_pairing_adjoint}). By the Choi trace pairing through the
    trace-pairing adjoint (\ref{thm:trace_choiMatrix_omega_quadratic_adjoint}),
    $\bra{\Omega}(T\otimes\Id)(\rho)\ket{\Omega}=\tr(W\rho)$, whose real part is
    negative. A positive semidefinite matrix has non-negative quadratic forms, so
    $(T\otimes\Id)(\rho)$ is not positive semidefinite.
\end{proof}

\section{Positive Schwarz maps outside complete positivity}

The following example shows that the Schwarz inequality does not force complete
positivity.

\subsection{A positive Schwarz map that is not completely positive}

\begin{example}[Transpose-trace map on $\MN{2}$]\label{def:wolf_example53}
    \lean{wolfExample53}
    \leanok
    Consider the linear map $T : \MN{2} \to \MN{2}$ defined by
    \begin{align}
        T(A) & =\frac{1}{2} A^T+\frac{1}{4}\tr(A)\Id.
        \notag
    \end{align}
    This is the map from \cite[Example~5.3]{Wolf2012Quantum}.
\end{example}

\begin{theorem}[The transpose-trace map is positive]\label{thm:wolf_example53_positive}
    \lean{wolfExample53_isPositive}
    \leanok
    \uses{def:wolf_example53, def:positive_map}
    The map from Example~\ref{def:wolf_example53} is positive.
\end{theorem}

\begin{proof}
    \leanok
    It is a non-negative linear combination of the transpose map and the map
    $A \mapsto \tr(A) \Id$, and both of these maps send positive
    semidefinite matrices to positive semidefinite matrices.
\end{proof}

\begin{theorem}[The transpose-trace map satisfies the Schwarz inequality]\label{thm:wolf_example53_schwarz}
    \lean{wolfExample53_satisfies_schwarz}
    \leanok
    \uses{def:wolf_example53}
    For every $A \in \MN{2}$,
    $T(A^\dagger A)-T(A^\dagger)T(A)\ge 0$, where $T$ is the map from
    Example~\ref{def:wolf_example53}.
\end{theorem}

\begin{proof}
    \leanok
    Write the Schwarz gap as $F(A)$.
    One checks that $F(A + \lambda \Id) = F(A)$ for every scalar $\lambda$,
    so it is enough to treat the case $\tr(A) = 0$.
    In that case a direct $2 \times 2$ computation shows that
    $F(A)$ dominates $\frac{1}{2} (A^\dagger A)^T$, hence is positive semidefinite.
\end{proof}

\begin{theorem}[The transpose-trace map is not completely positive]\label{thm:wolf_example53_not_cp}
    \lean{wolfExample53_not_cp}
    \leanok
    \uses{def:wolf_example53, def:cp_map}
    The map from Example~\ref{def:wolf_example53} is not completely positive.
\end{theorem}

\begin{proof}
    \leanok
    Compute the Choi matrix of $T$ and test it on the antisymmetric vector
    $\ket{01} - \ket{10}$.
    The resulting expectation value is negative, so the Choi matrix is not
    positive semidefinite.
\end{proof}
